%=========== ΛΥΣΗ - 59ο ΘΕΜΑ Δ ===========
\begin{thema}{Δ}
\begin{erwthma}
\item Για κάθε $x\in\mathbb{R}$ έχουμε
\begin{align*}
\left(\frac{1}{2}F^2(x)\right)'
&=F(x)F'(x)\\
&=F(x)f(x)\\
&=x.
\end{align*}
Επομένως
\[
\frac{1}{2}F^2(x)=\frac{x^2}{2}+c,
\]
όπου $c\in\mathbb{R}$. Για $x=0$ και με δεδομένο ότι $F(0)=1$, παίρνουμε
\[
\frac12=c.
\]
Άρα
\[
F^2(x)=x^2+1.
\]

Επειδή $x^2+1>0$, η $F$ δεν μηδενίζεται στο $\mathbb{R}$. Η $F$ είναι συνεχής και $F(0)=1>0$, οπότε διατηρεί θετικό πρόσημο. Συνεπώς
\[
{F(x)=\sqrt{x^2+1}},\quad x\in\mathbb{R}.
\]

\item Επειδή $f(x)=F'(x)$, έχουμε
\[
{f(x)=\frac{x}{\sqrt{x^2+1}}},\quad x\in\mathbb{R}.
\]

Για κάθε $x\in\mathbb{R}$ είναι
\[
f(-x)=\frac{-x}{\sqrt{x^2+1}}=-f(x).
\]
Άρα η $f$ είναι \textbf{περιττή}.

Επιπλέον,
\[
\lim_{x\to+\infty}\frac{x}{\sqrt{x^2+1}}
=\lim_{x\to+\infty}
\frac{1}{\sqrt{1+\frac{1}{x^2}}}
=1.
\]
Άρα η ευθεία
\[
{y=1}
\]
είναι οριζόντια ασύμπτωτη της $C_f$ στο $+\infty$.

Ομοίως,
\[
\lim_{x\to-\infty}\frac{x}{\sqrt{x^2+1}}
=-\lim_{x\to-\infty}
\frac{1}{\sqrt{1+\frac{1}{x^2}}}
=-1.
\]
Άρα η ευθεία
\[
{y=-1}
\]
είναι οριζόντια ασύμπτωτη της $C_f$ στο $-\infty$.

\item Για $x\to+\infty$ έχουμε
\begin{align*}
\lim_{x\to+\infty}\left[F(x)-x\right]
&=\lim_{x\to+\infty}
\left(\sqrt{x^2+1}-x\right)\\
&=\lim_{x\to+\infty}
\frac{1}{\sqrt{x^2+1}+x}\\
&=0.
\end{align*}
Επομένως η ευθεία
\[
{y=x}
\]
είναι πλάγια ασύμπτωτη της $C_F$ στο $+\infty$.

Για $x\to-\infty$ έχουμε
\begin{align*}
\lim_{x\to-\infty}\left[F(x)+x\right]
&=\lim_{x\to-\infty}
\left(\sqrt{x^2+1}+x\right)\\
&=\lim_{x\to-\infty}
\frac{1}{\sqrt{x^2+1}-x}\\
&=0.
\end{align*}
Επομένως η ευθεία
\[
{y=-x}
\]
είναι πλάγια ασύμπτωτη της $C_F$ στο $-\infty$.

\item Από τη δοσμένη σχέση έχουμε
\[
F(x)f(x)=x.
\]
Άρα
\[
xF(x)f(x)=x^2.
\]
Επομένως
\begin{align*}
\int_0^1xF(x)f(x)\d x
&=\int_0^1x^2\d x\\
&=\left[\frac{x^3}{3}\right]_0^1\\
&={\frac13}.
\end{align*}
\end{erwthma}
\end{thema}
%=========== ΤΕΛΟΣ ΛΥΣΗΣ - 59ο ΘΕΜΑ Δ ===========
